% Vietnamese translation for OI Public Library
% Source-File: IOI中国国家候选队论文/2004/朱晨光.ppt
\documentclass[11pt]{article}
\usepackage{oipl-vn}

\newcommand{\Oh}{\mathcal{O}}

\title{Tối ưu hóa quy hoạch động qua bài Trứng đại bàng\\{\large Ghi chú theo slide}}
\author{Zhu Chenguang}
\date{IOI 2004}

\begin{document}
\maketitle

\origtitle{Optimize, and optimize again!}
\sourcepath{IOI 2004 / Zhu Chenguang.ppt}

\section{Bài toán}

Các slide dùng bài thả trứng làm trục phân tích. Có \(M\) quả trứng và tòa
nhà \(N\) tầng, với \(M,N\le 1000\). Cần xác định ngưỡng bền \(E\) trong số
lần thả ít nhất ở trường hợp xấu nhất. Nếu \(M=1,N=10\), đáp án là \(10\),
vì chỉ có thể thử tuần tự từ thấp lên cao.

\section{Các bước tối ưu}

Trạng thái trực tiếp thường viết \(f(i,j)\): số lần thả ít nhất khi còn
\(i\) quả trứng và \(j\) tầng. Chọn một tầng thử tạo ra hai nhánh: trứng vỡ
thì còn ít trứng hơn và xét phần dưới; trứng không vỡ thì giữ nguyên số trứng
và xét phần trên. Công thức tự nhiên vì thế có dạng lấy cực tiểu của cực đại
hai phía.

Slide nhắc độ phức tạp ban đầu \(\Oh(n^3)\). Bản dịch PDF trình bày quá
trình hạ dần chi phí bằng các quan sát đơn điệu, đổi trạng thái sang số tầng
có thể kiểm soát trong số lần thả cho trước, rồi dùng quan hệ truy hồi
kiểu tổ hợp. Điểm đáng học là mỗi lần tối ưu đều xuất phát từ một ý nghĩa mới
của trạng thái, không chỉ từ việc viết lại công thức.

\section{Bài học thuật toán}

Với quy hoạch động, công thức đầu tiên thường chỉ xác nhận rằng bài toán có
cấu trúc con tối ưu. Muốn đạt hiệu năng tốt hơn, cần xét xem giá trị đang tối
ưu có thể đảo vai trò với tài nguyên hay không: thay vì hỏi ``cần bao nhiêu
lần thả cho \(N\) tầng'', hãy hỏi ``với \(t\) lần thả có thể phân biệt bao
nhiêu tầng''. Cách đổi góc nhìn này là điểm chính của bài trình bày.

\end{document}
